\chapter*{Bibliographical Studies}
\markboth{Bibliographical Studies}{Bibliographical Studies}

\section{UWB}

\subsection{Definitions} 
The Ultra-Wideband (UWB) is a technology using radio frequency to exchange data created in 1880.
The communication is not the only usage, it can also be used to locate, it's only  possible by using different modules, those algorithms and technique will be treated in the next part (\ref{UWB_Algo}).
In majority of cases, UWB is used for interior purpose, otherwise, it can be used for exterior purpose, but it's less shown in scientific paper.


\subsection{Location algorithm}\label{UWB_Algo}
In order to use the UWB as a location sensor. We have to use specific algorithm, in the followings of the subsection, we will explain and shown 2 diferents location algorithms. Firstly, the two way ranging (TWR) and secondly the double sided two way ranging (DS-TWR). And finally, we will explain TDOA, the most advanced algorithms.

\subsubsection{Two way ranging (TWR)}
The Two Way Ranging (TWR) algorithm is the easiest and the fastest to implement in a UWB system.
In order to explain this algorithm,  we will use 2 modules called $UWB_1$ and $UWB_2$.

\begin{figure}[!ht]
\centering
\resizebox{0.7\textwidth}{!}{%
\begin{circuitikz}
\tikzstyle{every node}=[font=\fontsize{18.2pt}{23.7pt}\selectfont]
\draw  (5,9) circle (1.25cm);
\draw  (15,9) circle (1.25cm);
\draw [-{Stealth[scale=1.5]}, ] (6.25,10.25) -- (13.75,10.25);
\draw [-{Stealth[scale=1.5]}, ] (13.75,7.75) -- (6.25,7.75);
\node [font=\fontsize{18.2pt}{23.7pt}\selectfont, inner xsep=0.080cm, inner ysep=0.085cm, rounded corners=0.020cm] at (14.875,9) {UWB};
\node [font=\fontsize{18.2pt}{23.7pt}\selectfont, inner xsep=0.080cm, inner ysep=0.085cm, rounded corners=0.020cm] at (4.875,9) {UWB};
\node [font=\fontsize{18.2pt}{23.7pt}\selectfont, inner xsep=0.080cm, inner ysep=0.085cm, rounded corners=0.020cm] at (10,10.75) {POLL};
\node [font=\fontsize{15.9pt}{20.7pt}\selectfont, inner xsep=0.080cm, inner ysep=0.085cm, rounded corners=0.020cm] at (5.75,8.875) {1};
\node [font=\fontsize{15.9pt}{20.7pt}\selectfont, inner xsep=0.080cm, inner ysep=0.085cm, rounded corners=0.020cm] at (15.75,8.875) {2};
\node [font=\fontsize{18.2pt}{23.7pt}\selectfont, inner xsep=0.080cm, inner ysep=0.085cm, rounded corners=0.020cm] at (10,7.25) {RESP};
\node [font=\fontsize{18.2pt}{23.7pt}\selectfont, inner xsep=0.080cm, inner ysep=0.085cm, rounded corners=0.020cm] at (6.125,10.875) {$t_{x1}$};
\node [font=\fontsize{18.2pt}{23.7pt}\selectfont, inner xsep=0.080cm, inner ysep=0.085cm, rounded corners=0.020cm] at (6.125,7.25) {$r_{x1}$};
\node [font=\fontsize{18.2pt}{23.7pt}\selectfont, inner xsep=0.080cm, inner ysep=0.085cm, rounded corners=0.020cm] at (13.875,10.875) {$r_{x2}$};
\node [font=\fontsize{18.2pt}{23.7pt}\selectfont, inner xsep=0.080cm, inner ysep=0.085cm, rounded corners=0.020cm] at (13.875,7.25) {$t_{x2}$};
\end{circuitikz}
}%
\caption{TWR explanation diagram}
\label{fig:TWR}
\end{figure}
\noindent The algorithm works in 4 steps :  
\begin{enumerate}
    \item $UWB_1$ sends a \textbf{POLL} to $UWB_2$. The time of emission $t_{x1}$ is memorized by $UWB_1$.
    \item $UWB_2$ receives the \textbf{POLL} and memorizes the time of arrival $r_{x2}$.
    \item $UWB_2$ processes the \textbf{POLL}, and precisely plans the time of emission $t_{x2}$ for the \textbf{RESP}. It contains ($t_{x2},r_{x2}$).
    \item $UWB_1$ receives the \textbf{RESP} and memorizes the $r_{x1}$.
\end{enumerate}
With all those parameters, we can deduce those variables :
\begin{itemize}
    \item Total time : $t_{tot} = r_{x1} - t_{x1}$
    \item Processing time of $UWB_2 : t_{proc} = r_{x2} - t_{x2}$
    \item Travel time : $t_{travel} = t_{tot} - t_{proc}$
\end{itemize}
The signal travels back and forth so it did twice the distance ($d$), and travel speed of a radio signal is equivalent to light speed ($c$). \\ 
So, we can deduce this equation : 
\begin{equation} \label{TWR_eq}
\boxed{\frac{t_{travel}\times c}{2}=\frac{(t_{tot} - t_{proc}) \times c}{2}= \frac{((r_{x1} - t_{x1}) - (r_{x2} - t_{x2})) \times c}{2}}
\end{equation}
We can also determine the induced error of this algorithm.\\
After calculation (describe in calculation part at page~\pageref{TWR-Error-calculation}), we obtain this equation: 
\begin{equation}
\epsilon \approx \frac{t_ {proc}(e_a - e_b)}{2}
\end{equation}
As we can observe, the Time of processing is directly link to the error. The $t_{proc}$ is around 3,000,000 nanoseconds so, it's inducing a high error. 
 
\subsubsection{Double Sided Two way ranging (DS-TWR)}
The DS-TWR has for purpose to reduce the error introduce by the $t_{proc}$. This algorithm also use $UWB_1$ and $UWB_2$.

\begin{figure}[!ht]
\centering
\resizebox{0.7\textwidth}{!}{%
\begin{circuitikz}
\tikzstyle{every node}=[font=\fontsize{18.2pt}{23.7pt}\selectfont]
\draw  (5,9) circle (1.25cm);
\draw  (15,9) circle (1.25cm);
\draw [-{Stealth[scale=1.5]}, ] (13.75,7.75) -- (6.25,7.75);
\node [font=\fontsize{18.2pt}{23.7pt}\selectfont, inner xsep=0.080cm, inner ysep=0.085cm, rounded corners=0.020cm] at (14.875,9) {UWB};
\node [font=\fontsize{18.2pt}{23.7pt}\selectfont, inner xsep=0.080cm, inner ysep=0.085cm, rounded corners=0.020cm] at (4.875,9) {UWB};
\node [font=\fontsize{18.2pt}{23.7pt}\selectfont, inner xsep=0.080cm, inner ysep=0.085cm, rounded corners=0.020cm] at (10,9.5) {POLL};
\node [font=\fontsize{15.9pt}{20.7pt}\selectfont, inner xsep=0.080cm, inner ysep=0.085cm, rounded corners=0.020cm] at (5.75,8.875) {1};
\node [font=\fontsize{15.9pt}{20.7pt}\selectfont, inner xsep=0.080cm, inner ysep=0.085cm, rounded corners=0.020cm] at (15.75,8.875) {2};
\node [font=\fontsize{18.2pt}{23.7pt}\selectfont, inner xsep=0.080cm, inner ysep=0.085cm, rounded corners=0.020cm] at (10,7.25) {RESP};
\node [font=\fontsize{18.2pt}{23.7pt}\selectfont, inner xsep=0.080cm, inner ysep=0.085cm, rounded corners=0.020cm] at (6.125,10.875) {$t_{x1}$};
\node [font=\fontsize{18.2pt}{23.7pt}\selectfont, inner xsep=0.080cm, inner ysep=0.085cm, rounded corners=0.020cm] at (6.125,7.25) {$r_{x1}$};
\node [font=\fontsize{18.2pt}{23.7pt}\selectfont, inner xsep=0.080cm, inner ysep=0.085cm, rounded corners=0.020cm] at (13.875,10.875) {$r_{x2}$};
\node [font=\fontsize{18.2pt}{23.7pt}\selectfont, inner xsep=0.080cm, inner ysep=0.085cm, rounded corners=0.020cm] at (13.875,7.25) {$t_{x2}$};
\draw [{Stealth[scale=1.5]}-, ] (13.625,10.125) -- (6.125,10.125);
\draw [{Stealth[scale=1.5]}-, ] (13.5,11.75) -- (6,11.75);
\node [font=\fontsize{18.2pt}{23.7pt}\selectfont, inner xsep=0.080cm, inner ysep=0.085cm, rounded corners=0.020cm] at (9.625,11.25) {FINAL};
\node [font=\fontsize{18.2pt}{23.7pt}\selectfont, inner xsep=0.080cm, inner ysep=0.085cm, rounded corners=0.020cm] at (6.125,12.5) {$t_{x12}$};
\node [font=\fontsize{18.2pt}{23.7pt}\selectfont, inner xsep=0.080cm, inner ysep=0.085cm, rounded corners=0.020cm] at (13.75,12.5) {$r_{x22}$};
\end{circuitikz}
}%
\caption{DS-TWR explanation diagram}
\label{fig:DS-TWR}
\end{figure}

\subsubsection{TDOA}

\newpage
\section{Measure treatment}
In most cases, date measured by sensors are unclear. It could be a sensor offset, it can also come from crosstalk. Or others, exterior interaction. In order, to obtain cleans measurements, and to be able to used them for controlling or any other purpose. We have to develop some filter, in this section, we will explain three filters used for cleaning UWB signal, Firstly, we will explain the polynomial filter, and secondly, the Kalman filter, then finally the extended Kalman filter
\subsection{Polynomial Filter}


\subsection{Kalman Filter}

\subsection{Extended Kalman Filter}

\newpage
